\documentclass[12pt]{article}
\usepackage[a4paper,margin=1in]{geometry}
\usepackage{amsmath,amssymb}
\begin{document}

%---------------------------------------
% 1. Parametric (Mathieu) Equation
%---------------------------------------
\section*{Parametric (Mathieu) Equation and Resonance}
\begin{equation*}
  \ddot{x} \;+\;\omega_{0}^{2}\bigl[\,1 + \epsilon\,\cos(\gamma\,t)\bigr]\,x \;=\; 0.
\end{equation*}
A trial solution of the form
\[
x(t + nT) \;=\; \mu^{n}\,x(t)
\]
shows that if $|\mu| \neq 1$, the solution grows or decays, indicating \emph{parametric resonance}.

Sometimes this is also written in the form
\[
x'' + \bigl[\omega_{0}^{2} + \epsilon \cos(\gamma t)\bigr]x = 0,
\]
known as the Mathieu equation when $\gamma$ is the driving frequency.

A related Wronskian-type invariant for second-order linear ODEs:
\[
x_{2}\,\frac{d x_{1}}{dt} \;-\; x_{1}\,\frac{d x_{2}}{dt} \;=\;\text{constant}.
\]

%---------------------------------------
% 2. Scattering Cross Section
%---------------------------------------
\section*{Scattering Cross Section}
For scattering with impact parameter $\rho$, scattering angle $\theta$, asymptotic velocity $v_\infty$, and reduced mass $\mu$, the differential cross section is:
\begin{equation*}
    \frac{d\sigma}{d\Omega}
    \;=\;
    \frac{\rho(\theta)}{\sin \theta}
    \left|\frac{d\rho}{d\theta}\right|.
\end{equation*}
The total cross section is:
\begin{equation*}
    \sigma_{T}
    \;=\;
    2\pi \int_{0}^{\pi}
    \rho(\theta)
    \left|\frac{d\rho}{d\theta}\right|
    \,d\theta.
\end{equation*}

We identify
\[
E \;=\;\frac12\,\mu\,v_{\infty}^{2},
\quad
\ell \;=\;\mu\,\rho\,v_{\infty},
\]
and if $r_{\min}$ is the distance of closest approach in the scattering, then
\[
\rho^{2}
\;=\;
r_{\min}^{2}\Bigl(1 \;-\;\frac{V(r_{\min})}{E}\Bigr).
\]

%---------------------------------------
% 3. Symmetric Top Energies & Frequencies
%---------------------------------------
\section*{Symmetric Top Energies and Frequencies}
For a symmetric top with principal moments $I_{1} = I_{2} \neq I_{3}$ and total angular momentum $J$
\begin{equation*}
    T
    \;=\;
    \frac{J^{2}}{2}\,
    \Bigl(\frac{1}{I_{1}}
          \;+\;
          \cos^{2}\theta\,
          \bigl[\tfrac{1}{I_{3}}-\tfrac{1}{I_{1}}\bigr]
    \Bigr).
\end{equation*}
In the absence of external torque (or with the appropriate constraints),
\begin{equation*}
    \omega_{r}
    \;=\;
    \omega_{3}
    \;-\;
    \frac{J}{I_{1}}\cos\theta,
    \qquad
    \omega_{p}
    \;=\;
    \frac{J}{I_{1}}.
\end{equation*}

%---------------------------------------
% 4. Two-Timescale (Slow + Fast) Splitting
%---------------------------------------
\section*{Two-Timescale Splitting and Effective Potential}
In certain problems, one can split a coordinate $\phi$ into a slowly varying part $\Phi$ plus a fast oscillation $\eta$, i.e.\ $\phi(t) = \Phi(t) + \eta(t).$
If
\[
\phi'' = -\frac{\partial V}{\partial \phi} + f(\phi,t),
\]
then, by averaging over the fast oscillations, one defines an \emph{effective} potential:
\begin{equation*}
    V_{\text{eff}}(\Phi)
    \;=\;
    V(\Phi)
    \;+\;
    \tfrac{1}{2}\,\bigl\langle \dot{\eta}^{2}\bigr\rangle.
\end{equation*}
The equilibrium points of $V_{\text{eff}}$ give approximate slow solutions (sometimes called a “secular” approximation).

%---------------------------------------
% 5. Additional Canonical Identities
%---------------------------------------
\section*{Additional Canonical Transformations and Identities}
If $F$ is a generating function, then
\begin{equation*}
    H' \;=\; H \;+\; \frac{\partial F}{\partial t}.
\end{equation*}
A small canonical transformation generated by $G(p,q)$ changes any phase-space function $u$ by
\begin{equation*}
    \delta u
    \;=\;
    \epsilon\,\{\,u,\,G\},
\end{equation*}
where $\{\cdot,\cdot\}$ is the Poisson bracket.  If $A$ is a constant of motion ($\{H,A\}=0$), then $A$ generates a continuous symmetry of the Hamiltonian.

\end{document}
